
\def\PUDcodigo{EME103}

\def\PUDcodacad{04506.54}

\def\PUDdisciplina{Transportadores Industriais}

\def\PUDsemestre{Opt}

\def\PUDhoras{80}

%NOVOS ---------------------------------------------------
\def\PUDhorasTeoria{80}
\def\PUDhorasPratica{0}

\def\PUDinterECA{---}

\def\PUDatuaisECA{---}

\def\PUDrecursos{Projetor multimídia; Quadro branco e pincel.}

%----------------------------------------------------------

\def\PUDcreditos{4}

\def\PUDprerequisito{ \hyperref[PUD:EME119]{EME119} }

\def\PUDpreacad{ \hyperref[PUD:EME119]{04506.35} }

\def\PUDobjetivos{Conhecer as classes de cargas industriais;  Conhecer os principais tipos de transporte de cargas industriais;  Dimensionar os principais tipos de transportadores industriais.}

\def\PUDementa{Introdução à Movimentação de Materiais;  Especificação
de Transportadores Industriais;  Dimensionamenton de Transportadores Industriais;  Considerações
sobre Transporte Industrial;  Transportadores Contínuos;  Componentes das Máquinas de Elevação.}

\def\PUDconteudo{ & Introdução à movimentação de materiais.  Elementos fundamentais em máquinas de elevação e transporte.  
\\ 
 & Especificação de transportadores industriais: Natureza da carga, parâmetros de transporte. 
\\ 
 & Dimensionamento de transportadores industriais.  Posicionadores, projeto e fabricação de transportadores industriais.  Normas de segurança. 
\\ 
 & Considerações sobre transporte industrial.  Equipamentos para transporte, transferência, condução e elevação.  
\\ 
 & Transportadores contínuos, correias, capacidade do transportador, sistemas de acionamento, roletes.  Classificação das máquinas de elevação.  
\\ 
 & Componentes das máquinas de elevação.  Dispositivos de apanhar carga.  Mecanismos de elevação.  Freios.  Mecanismos de translação.  Pontes rolantes. }

\def\PUDmetodo{Aulas expositórias que podem ser teóricas e/ou práticas, onde as práticas no laboratório serão marcadas no decorrer da disciplina.}

\def\PUDavalia{A avaliação será feita com aplicação de provas teóricas e/ou práticas, além da possibilidade de inclusão de trabalhos e seminários no decorrer da disciplina.}

\def\PUDbibbasica{ & \href{www.biblioteca.ifce.edu.br/index.asp?codigo_sophia=40282}{Norton}, Robert L.  Projeto de máquinas: uma abordagem integrada.  4º ed.  São Paulo - SP.  Editora BOOKMAN.  2013.  1028p.  ISBN: 9788582600221. 
\\ 
 & \href{www.biblioteca.ifce.edu.br/index.asp?codigo_sophia=23902}{Collins}, Jack A.  Projeto mecânico de elementos de máquinas: uma perspectiva de prevenção da falha.  1ª edição.  Rio de Janeiro - RJ.  Editora LTC.  2006.  740p.  ISBN: 8521614756. 
\\ 
 &  \href{www.biblioteca.ifce.edu.br/index.asp?codigo_sophia=63128}{Juvinall}, Robert C.  Fundamentos do Projeto de Componentes de Máquinas.  5º ed.  Rio de Janeiro - RJ.  Editora LTC.  2016.  ISBN: 9788521630098. }

\def\PUDbibcomplementar{ &  \href{www.biblioteca.ifce.edu.br/index.asp?codigo_sophia=62907}{Macintyre}, Archibald Joseph.  Equipamentos Industriais e de Processo.  1ª edição.  Rio de Janeiro - RJ.  Editora LTC.  2016.  277p.  ISBN: 9788521611073.
\\
 &  \href{www.biblioteca.ifce.edu.br/index.asp?codigo_sophia=63248}{Groover}, Mikell P.  Introdução aos processos de fabricação.  Rio de Janeiro, RJ : LTC, 2014.  737p.  ISBN: 9788521625193.
%  RUDENKO, N.;  Máquinas de Elevação e Transporte.  5ª edição.  Rio de Janeiro.  Editora LTC.  1976.  425p.  ISBN:
\\ 
 &  \href{www.biblioteca.ifce.edu.br/index.asp?codigo_sophia=65223}{Bateman}, Robert E.  et al.  Simulação de sistemas: aprimorando processo de logística, serviços e manufatura.  Rio de Janeiro, RJ : Elsevier, 2013.  161p.  ISBN: 9788535271621.
% SHIGLEY, J.  E.;  MISCHKE, C.  R.;  BUDYNAS, R. G.  Projeto de Engenharia Mecânica.  7ª edição, Porto Alegre Editora Bookman, 2007 960p ISBN: 9788536305622.
\\
 &  \href{www.biblioteca.ifce.edu.br/index.asp?codigo_sophia=39529}{Fogliatti}, Maria Cristina.  Avaliação de impactos ambientais: aplicação aos sistemas de transporte.  Rio de Janeiro, RJ: Interciência, 2004.  249p.  ISBN: 8571931089.
\\
 &  \href{www.biblioteca.ifce.edu.br/index.asp?codigo_sophia=40619}{Livi}, Celso Pohlmann.  Fundamentos de fenômenos de transporte: um texto para cursos básicos.  2º ed.  Rio de Janeiro, RJ: LTC, 2012.  237p.  ISBN: 9788521620570. }

\def\PUDautor{Venceslau Xavier de Lima Filho}

\def\PUDdata{2013-06-17}

\def\PUDrevisao{1}

\def\PUDrevisor{Samuel Vieira Dias}

\def\PUDdatarevisao{2017-04-30}

\PUDcorpo
